\paragraph{Credit and coordinate comparison.}
The cuspidal Weierstrass self-product and its projective small resolutions
are the construction of Namikawa described in
\cite[Proposition~4, pp.~10--11]{Rossi2017}; the ambient Hodge numbers
$h^{1,1}=21$ and $h^{2,1}=3$ are also recorded in
\cite[Remark~5, p.~20]{Rossi2017}.
The coordinate exchange $(x_{\rm Rossi},y_{\rm Rossi})=(y,x)$
and the coefficient $B(t)=t^6-1$ identify the underlying Weierstrass
surface and its fibre self-product with the equations below.
We retain the explicit two blowup ideals, the full support, and its
integral collapse and normalization maps in the proof.

\begin{proposition}[A simply connected irreducible rational-normalized fibre]\label{cdp:st:C.2.UN.8}
There is a smooth projective threefold $Y$ with $K_Y\simeq\OO_Y$
and a connected-fibre map $f:Y\to\PP^1$ whose general fibre is a
product of elliptic curves. It has six singular fibres. Each is
reduced and irreducible, has rational normalization resolution, and
satisfies
\[
 \begin{gathered}
 \pi_1(S)=1,\qquad e(S)=6,\\
 (H^i(S,\Z))_{i=0}^4=(\Z,0,\Z^4,0,\Z),\qquad
 (h^i(S,\OO_S))_{i=0}^2=(1,2,1).
 \end{gathered}
 \tag{C.CP.1}\label{cdp:eq:C.CP.1}
\]
The ambient cohomology is
\[
 \begin{gathered}
 (b_i(Y))_{i=0}^6=(1,0,21,8,21,0,1),\qquad e(Y)=36,\\
 (h^i(Y,\OO_Y))_{i=0}^3=(1,0,0,1),\qquad
 h^{1,1}(Y)=21,\quad h^{2,1}(Y)=3,
 \end{gathered}
 \tag{C.CP.2}\label{cdp:eq:C.CP.2}
\]
where the Betti numbers use rational coefficients. In particular
the nonzero local character used in \ref{cdp:st:C.2.MH.2} need not exist after
the global CDP hypotheses are removed, and its Euler conclusion
$e(S)=2$ does not extend to every irreducible rational-normalized
torus degeneration.
\end{proposition}
\Uses{the hypotheses, notation and definitions stated here; \cite[Proposition~4, pp.~10--11, and Remark~5, p.~20]{Rossi2017}; \cref{cdp:st:C.2.MH.2}; \cref{cdp:eq:C.CP.5}; \cref{cdp:eq:C.CP.7}; \cref{cdp:eq:C.CP.6}; \cref{cdp:eq:C.CP.4}; \cref{cdp:eq:C.CP.10}}

\begin{proof}
Let $B=\PP^1$, let $L=\OO_B(1)$, and take the Weierstrass
surface $g:R\to B$ with zero section and equation
\[
                  y^2z=x^3+p(t)z^3,
             \qquad p(t)=t^6-1.                 \tag{C.CP.3}\label{cdp:eq:C.CP.3}
\]
Here $x,y$ have the usual transition weights $L^2,L^3$;
thus $p$ is a section of $L^6$. The projective Weierstrass
bundle supplies the compactification. At infinity the transformed
coefficient is $1-t'^6$, so that fibre is smooth. The six
zeros of $p$ are simple. The only singular fibres are the six
cuspidal cubics, and the surface is smooth: in the affine
Weierstrass chart a possible singular point has $x=y=0$, while
the derivative $p'(t)$ there is nonzero. The zero-section chart
is smooth directly from its equation. Adjunction in the
Weierstrass bundle gives
\[
                         K_R=g^*\OO_B(-1).
 \tag{C.CP.4}\label{cdp:eq:C.CP.4}
\]
This formula can also be read from the local canonical form
$dt\wedge dx/y$ and the transition weights just specified.

Form the projective fibre product $X=R\times_B R$. Its only
singularities are the six points where both factors are cusp
points. At one of them use the base coordinate $\tau=p(t)$,
and write
\[
 y_1^2-x_1^3=\tau=y_2^2-x_2^3.
\]
Fix a primitive cube root $\rho$. With
\[
 \begin{gathered}
 u=y_1-y_2,\quad v=y_1+y_2,\quad
 a=x_1-x_2,\quad b=x_1-\rho x_2,\\
 c=x_1-\rho^2x_2=\alpha a+\beta b,
 \qquad \alpha=-\rho,\quad\beta=-\rho^2,
 \end{gathered}
\]
the exact threefold equation is
\[
                         uv=abc.                \tag{C.CP.5}\label{cdp:eq:C.CP.5}
\]
Both $\alpha$ and $\beta$ are nonzero. We resolve \eqref{cdp:eq:C.CP.5}
by two specified blowups. First blow up the ideal $(u,a)$.
In its $u$-chart, $a=ur$ and
\[
                  v=rb(\alpha ur+\beta b),
 \tag{C.CP.6}\label{cdp:eq:C.CP.6}
\]
so $(u,r,b)$ are smooth coordinates. In its $a$-chart put
$u=as$. The equation is $sv=b(\alpha a+\beta b)$; its
only singular point is the origin, an ordinary threefold
node. The strict transform of $(u,b)$ has ideal $(s,b)$
there. Blowing up that ideal gives the two smooth charts
\[
 \begin{array}{c|c}
 \text{coordinates}&\text{blowdown formulas}\\ \hline
 (a,s,h)&b=sh,\quad u=as,\quad
                    v=h(\alpha a+\beta sh),\\
 (b,j,v)&s=bj,\quad a=(jv-\beta b)/\alpha,\quad u=abj.
 \end{array}                                    \tag{C.CP.7}\label{cdp:eq:C.CP.7}
\]
The inverse image of the original singular point is a chain
of two projective lines. In the first chart of \eqref{cdp:eq:C.CP.7}
the old and new curves have equations $(a,h)=(0,0)$
and $(a,s)=(0,0)$; they meet once transversely. Charts
\eqref{cdp:eq:C.CP.6}--\eqref{cdp:eq:C.CP.7} complete each to a projective line.
No exceptional divisor is introduced. Outside the origin
the respective ideals are Cartier, so these maps are
isomorphisms there.

The blowups globalize without an analytic gluing choice.
The ideal $(u,a)$ is the ideal of the Weil diagonal
$R\hookrightarrow X$. The ideal $(u,b)$ is the Weil
graph of the global automorphism
$(x,y,t)\mapsto(\rho x,y,t)$ of $R$; after the first
blowup take its strict transform. Each of these divisors
is Cartier off the six specified singular points.
Their projective blowups therefore give a smooth
projective threefold $Y$ and a proper small map
$\pi:Y\to X$, with precisely the local charts above.
Set $f=g\pi$, using the common projection of the fibre
product. It agrees with the original product of elliptic
curves away from the six values. Its fibres are connected,
because the product fibres and every fibre of the two
blowdowns are connected.

The original central support at each of the six values
is the product $S_0=C\times C$ of a cuspidal cubic with
itself. It is reduced and irreducible. Since $\pi$ has
no exceptional divisors, the full fibre upstairs is its
single strict transform, still reduced. Its normalization
resolution is rational: the function field is unchanged
from that of the finite birational parametrization
$\PP^1\times\PP^1\to S_0$. This assertion is birational;
we do not discard the exceptional curves of the actual
normalization resolution.

Topologically a cuspidal cubic is $S^2$. Its finite
normalization is bijective, so compactness makes it a
homeomorphism. Hence $S_0$ is homeomorphic to
$S^2\times S^2$. The restriction $S\to S_0$ is an
isomorphism off the cusp-pair point and replaces that
point by the chain $E$ of two projective lines. The
quotient $S/E\to S_0$ is a continuous bijection from a
compact space to a Hausdorff space, hence a homeomorphism.
Choose a finite triangulation of $(S,E)$. Its cofibration
sequence identifies the reduced homology of $S/E$ with
the relative homology of $(S,E)$. Since
$H_2(E,\Z)=\Z^2$ and $E$ is simply connected, the
long exact sequence gives
\[
 \begin{gathered}
 H_4(S,\Z)=\Z,\qquad H_3(S,\Z)=H_1(S,\Z)=0,\\
 0\longrightarrow\Z^2\longrightarrow H_2(S,\Z)
                 \longrightarrow\Z^2\longrightarrow0.
 \end{gathered}                                  \tag{C.CP.8}\label{cdp:eq:C.CP.8}
\]
The last sequence splits as a sequence of abelian groups.
Van Kampen for the same cofibration shows
$\pi_1(S/E)=\pi_1(S)/\langle\!\langle\pi_1(E)\rangle\!\rangle$.
Thus $\pi_1(S)=1$. Universal coefficients now give the
integral cohomology in \eqref{cdp:eq:C.CP.1}, and its Euler number
is $1+4+1=6$.

The fibre product is a local complete intersection in
$R\times R$, cut out by the inverse image of the base
diagonal. Adjunction and \eqref{cdp:eq:C.CP.4} give
\[
 \omega_X=p_1^*K_R\otimes p_2^*K_R\otimes f^*T_B
                         \simeq\OO_X.
 \tag{C.CP.9}\label{cdp:eq:C.CP.9}
\]
Its generator pulls back to a nowhere-vanishing canonical
form on $Y$ away from the small exceptional curves.
Its coefficients and the coefficients of its inverse
extend across that codimension-two set by Hartogs'
theorem in a local canonical frame. Hence $K_Y=\OO_Y$.
Adjunction makes $\omega_S=\OO_S$. Proper flatness over
the smooth base preserves the holomorphic Euler
characteristic, so $\chi(\OO_S)=\chi(\OO_{E_t^2})=0$.
Connectedness and Serre duality then give
$(h^0,h^1,h^2)(\OO_S)=(1,2,1)$.

We include the ambient groups, with their comparison
maps. A type-II local elliptic monodromy has characteristic
polynomial $x^2-x+1$, so $T-I$ is an integral isomorphism.
This also follows directly from the cyclic sixth-root
trivialization of $y^2=x^3+\tau$. The smooth part of
the fibration has fundamental group generated by its
fibre lattice and base meridians. Its inclusion into
$Y$ is surjective on fundamental groups, by perturbing
loops off the central divisor. The pair of zero sections
extends over every special value and kills each base
meridian in $Y$. Conjugation by one such meridian
therefore kills the image of $T-I$ on each of the two
elliptic lattices. Since those two maps are unimodular,
the whole fibre lattice is killed. Thus $\pi_1(Y)=1$.
Projectivity and $K_Y=\OO_Y$ give
$(h^0,h^1,h^2,h^3)(\OO_Y)=(1,0,0,1)$ by Hodge theory
and Serre duality. The identical argument for $R$ gives
$b_1(R)=0$. Its six cuspidal fibres give $e(R)=12$,
hence $b_2(R)=10$. Moreover $K_R=g^*\OO(-1)$ makes
every positive plurigenus zero. The surface classification
therefore makes $R$ rational, so its Picard group has
rank ten.

Let $K=\C(t)$ and let $E_K$ be the generic elliptic
curve of $R$. Restriction of divisors gives a surjection
$\Pic(R)\to\Pic(E_K)$: close a generic divisor in the
smooth surface. Its kernel is generated, up to linear
equivalence, by the full fibre, because every fibre is
irreducible. The zero section splits the degree map on
$\Pic(E_K)$, whose degree-zero part is $E_K(K)$.
Consequently $E_K(K)$ has rank eight.

Apply the same divisor restriction to $Y$ and its
generic fibre $A_K=E_K\times E_K$. It is surjective
and its vertical kernel has rank one, since all six
singular fibres of $Y$ are irreducible. The degree-zero
Picard group of $A_K$ is its dual abelian variety's
$K$-points, which have rank sixteen under the product
principal polarization. Its N\'eron--Severi group has
rank four. To verify this last number explicitly, use
the two factors, the diagonal, and the graph of the
automorphism $\rho$. Their intersection matrix is
\[
 \begin{pmatrix}
 0&1&1&1\\1&0&1&1\\1&1&0&3\\1&1&3&0
 \end{pmatrix}.                                \tag{C.CP.10}\label{cdp:eq:C.CP.10}
\]
Its determinant is $-3$. The entry three is the degree of
$1-\rho$, whose kernel on the triangular elliptic
lattice has order three. Here is the exact comparison
for the upper bound on the generic N\'eron--Severi rank.
Over $K'=K(u)$ with $u^6=p(t)$, the algebraic
coordinate substitution $x=u^2X$, $y=u^3Y$ identifies
$E_{K'}$ with the constant elliptic curve
$E_0:\ Y^2=X^3+1$. Base change to an algebraic closure
identifies the geometric N\'eron--Severi group with that
of $(E_0\times E_0)_{\overline K}$. For the constant
complex abelian surface, the first Chern class embeds
this group into the integral $(1,1)$ classes and
$h^{1,1}(E_0^2)=4$. Algebraically closed extension
preserves the N\'eron--Severi group: the Picard scheme
and its identity component commute with field extension,
and each of its connected components over $\C$ remains
geometrically connected. The four classes
in \eqref{cdp:eq:C.CP.10} are already defined over $K$, so its rank
over $K$ is exactly four. Thus
$\operatorname{rank}\Pic(Y)=1+16+4=21$.
Since $h^1(\OO_Y)=h^2(\OO_Y)=0$, the exponential
sequence identifies this rank with $b_2(Y)$.

Euler localization now gives $e(Y)=6\cdot6=36$.
Poincar\'e duality and $b_1=0$ yield $b_3=8$;
the remaining Betti numbers and Hodge numbers follow
as displayed in \eqref{cdp:eq:C.CP.2}. These rational Betti assertions
make no assertion about integral torsion in degrees
not already computed.

Finally a sufficiently small compact fibre tube retracts
onto $S$. Its fundamental group is therefore trivial.
There is no nontrivial finite local system there, and
no nonzero integral degree-one class to restrict to
the nearby torus. The example consequently disproves
the proposed unrestricted replacement for the character
construction itself, with its actual comparison map.
The exact coordinate identities, intersection matrix,
and all displayed integer ranks are checked in
\path{code/root_cusp_product.py} and its saved output.
\end{proof}